\documentclass[12pt]{article}

% Match the earlier reports using locally available fonts.
\usepackage[T1]{fontenc}
\usepackage{times}
\usepackage{microtype}
\usepackage{amsmath}
\usepackage[parfill]{parskip}
\setlength{\parindent}{0pt}
\setlength{\parskip}{2ex plus 0.2ex minus 0.2ex}
\usepackage[letterpaper,margin=1in]{geometry}
\usepackage[hidelinks]{hyperref}

\begin{document}

\begin{flushleft}
\textbf{Week 9: Hierarchical Models, Robust Models, and Empirical Bayes} \\
Jaisel Singh (js6897) \\
November 11, 2026
\end{flushleft}

\vspace{0.1in}

What stood out to me in Gelman and Hill was that multilevel modeling gives the
regression coefficients their own model. Group intercepts might satisfy
$\alpha_j\sim\mathcal{N}(\mu,\tau^2)$, allowing groups to differ while sharing
information about their distribution. The radon example made partial pooling
concrete: an extreme county average based on two measurements should carry less
weight than an equally extreme average supported by many measurements. The
amount of shrinkage depends on the precision of the group's observations
relative to the variation between groups. I find this a useful way to think
about regularization: we stabilize uncertain estimates through an explicit
assumption about how the groups are related.

Efron's distinction between oracle Bayes and empirical Bayes clarified what we
gain and what we must estimate. The oracle supplies information about the
collection's true parameter values, giving a benchmark for estimation; empirical
Bayes learns the distribution used for inference from the observed collection.
This lets many related estimation problems inform each individual one.
However, the estimated prior is itself uncertain. Plugging it in as though it
were known omits a source of uncertainty, while a fully Bayesian hierarchy can
propagate uncertainty in its hyperparameters. I found Efron's separation of
overall estimation accuracy from inference about a particular case especially
helpful: improving average performance does not settle whether the comparison
group is relevant to that case.

The robustness theme makes me think about where unusual behavior occurs. A
heavy-tailed observation model can reduce an extreme measurement's influence
on a fitted location, but that is a different issue from an unusual
group. Pooling the latter toward an unsuitable population could obscure a real
difference. This connects Efron's relevance question to the hierarchy itself.
My remaining question is how to use model checks and subject knowledge to
decide whether an extreme group needs a different error model, additional
group-level predictors, or a different population to borrow information from.

\vspace{0.05in}
{\footnotesize
\textit{Readings:} Gelman and Hill (2007), \emph{Data Analysis Using Regression and Multilevel/Hierarchical Models}, Chapters 11--12;
Efron (2019), \href{https://doi.org/10.1214/18-STS674}{\emph{Bayes, Oracle Bayes and Empirical Bayes}}.
\par}

\end{document}

%%% Local Variables:
%%% mode: latex
%%% TeX-master: t
%%% End:
